% GENERATED FILE. Edit canonical lessons, then run npm run generate:books.
% canonical-source-hash: fnv1a64:048aee53ac5a6dcd
% canonical-lessons: RU-C415-otpravlyatsya, RU-C415-pribyvat, RU-C415-otkhodit, RU-C415-obyavlyat, RU-C415-uspevat

\chapter{The train is leaving}
\label{ch:ru-the-train-is-leaving}
\hlchaptermodality{415}

\begin{chapteropening}
\textbf{By the end of this chapter:} \emph{I can follow station announcements: setting off, arriving, a train pulling out, announcing and making it in time.}

Complete the last lesson of chapter 415: I can follow station announcements: setting off, arriving, a train pulling out, announcing and making it in time.
\end{chapteropening}

\section[\texorpdfstring{\ru{отправляться} (otpravlyátsya)}{отправляться (otpravlyátsya)}]{\ru{отправляться} (otpravlyátsya) — to set off, to depart}

\label{lesson:RU-C415-otpravlyatsya}



\begin{quote}
Before the new one: say the Russian for departure, then the Russian for arrival.
\end{quote}

\subsection*{You'll want to know: \ru{отправляться}}
\textbf{\ru{отправляться}} (\emph{otpravlyátsya}) — \textquotedblleft{}to set off, to depart\textquotedblright{}.

\begin{quote}
\ru{Поезд} \ru{отправляется} \ru{в} \ru{восемь} \ru{утра}.
\end{quote}

\begin{quote}
\emph{Póyezd otpravlyáyetsya v vósem utrá.}
\end{quote}

\begin{quote}
The train departs at eight in the morning.
\end{quote}

Imperfective; perfective \emph{\ru{отправиться}}. \emph{\ru{Отправлять}} is to send; with \emph{-\ru{ся}}, to set off oneself.

\begin{cousinweb}
\emph{\ru{отправляться}} is \emph{\ru{отправлять}} 'to send off' with \emph{-\ru{ся}} 'oneself'. \emph{\ru{Отправлять}} is built from \emph{\ru{от}-} 'off' and \emph{\ru{править}} 'to steer, to set right'; EtymDB lists the noun \emph{\ru{отправление}} 'departure'.
\end{cousinweb}

\begin{grammarlens}[title={What you've built}]
One more word for telling, explaining and describing at B1.
\end{grammarlens}

\subsection*{Your turn}
\begin{itemize}
\raggedright
  \item \emph{Say it:} \emph{otpravlyátsya}
  \item \emph{Say it:} \emph{otpravlyátsya}, once more
  \item \emph{Say it:} a sentence of your own with \emph{otpravlyátsya}
  \item \emph{Recall:} say the Russian for departure, then the Russian for arrival, then say \emph{otpravlyátsya} again
\end{itemize}

\begin{tcolorbox}[breakable,colback=teal!4,colframe=teal!35!black,title={Before you move on}]
What does \ru{отправляться} mean? (\textquotedblleft{}To set off, to depart\textquotedblright{}.) Say it once more.
\end{tcolorbox}
\section[\texorpdfstring{\ru{прибывать} (pribyvát)}{прибывать (pribyvát)}]{\ru{прибывать} (pribyvát) — to arrive (formal)}

\label{lesson:RU-C415-pribyvat}



\begin{quote}
Before the new one: say the Russian for arrival, then the Russian for to set off.
\end{quote}

\subsection*{You'll want to know: \ru{прибывать}}
\textbf{\ru{прибывать}} (\emph{pribyvát}) — \textquotedblleft{}to arrive (formal)\textquotedblright{}.

\begin{quote}
\ru{Самолёт} \ru{прибывает} \ru{в} \ru{аэропорт} \ru{в} \ru{полдень}.
\end{quote}

\begin{quote}
\emph{Samolyót pribyváyet v aeropórt v pólden.}
\end{quote}

\begin{quote}
The plane arrives at the airport at noon.
\end{quote}

Imperfective; perfective \emph{\ru{прибыть}} (\emph{\ru{прибуду}, \ru{прибудешь}}). Formal; in speech \emph{\ru{приезжать}}, \emph{\ru{прилетать}}.

\begin{cousinweb}
\emph{\ru{прибывать}} is built from \emph{\ru{при}-} 'to, near' and \emph{\ru{бывать}}, from \emph{\ru{быть}} 'to be'. EtymDB lists \emph{\ru{прибыть}} without a deeper source.
\end{cousinweb}

\begin{grammarlens}[title={What you've built}]
One more word for telling, explaining and describing at B1.
\end{grammarlens}

\subsection*{Your turn}
\begin{itemize}
\raggedright
  \item \emph{Say it:} \emph{pribyvát}
  \item \emph{Say it:} \emph{pribyvát}, once more
  \item \emph{Say it:} a sentence of your own with \emph{pribyvát}
  \item \emph{Recall:} say the Russian for arrival, then the Russian for to set off, then say \emph{pribyvát} again
\end{itemize}

\begin{tcolorbox}[breakable,colback=teal!4,colframe=teal!35!black,title={Before you move on}]
What does \ru{прибывать} mean? (\textquotedblleft{}To arrive (formal)\textquotedblright{}.) Say it once more.
\end{tcolorbox}
\section[\texorpdfstring{\ru{отходить} (otkhodít)}{отходить (otkhodít)}]{\ru{отходить} (otkhodít) — to pull out, to leave (of a train); to step away}

\label{lesson:RU-C415-otkhodit}



\begin{quote}
Before the new one: say the Russian for to set off, then the Russian for to arrive.
\end{quote}

\subsection*{You'll want to know: \ru{отходить}}
\textbf{\ru{отходить}} (\emph{otkhodít}) — \textquotedblleft{}to pull out, to leave (of a train); to step away\textquotedblright{}.

\begin{quote}
\ru{Наш} \ru{поезд} \ru{отходит} \ru{через} \ru{пять} \ru{минут}.
\end{quote}

\begin{quote}
\emph{Nash póyezd otkhódit chérez pyat minút.}
\end{quote}

\begin{quote}
Our train leaves in five minutes.
\end{quote}

Imperfective; the perfective is \emph{\ru{отойти}} (\emph{\ru{отойду}, \ru{отойдёшь}}). Present \emph{\ru{отхожу}, \ru{отходишь}}.

\begin{cousinweb}
\emph{\ru{отходить}} is built from \emph{\ru{от}-} 'away' and \emph{\ru{ходить}} 'to go'. EtymDB takes the root from Proto-Slavic \emph{xodŭ} 'motion'.
\end{cousinweb}

\begin{grammarlens}[title={What you've built}]
One more word for telling, explaining and describing at B1.
\end{grammarlens}

\subsection*{Your turn}
\begin{itemize}
\raggedright
  \item \emph{Say it:} \emph{otkhodít}
  \item \emph{Say it:} \emph{otkhodít}, once more
  \item \emph{Say it:} a sentence of your own with \emph{otkhodít}
  \item \emph{Recall:} say the Russian for to set off, then the Russian for to arrive, then say \emph{otkhodít} again
\end{itemize}

\begin{tcolorbox}[breakable,colback=teal!4,colframe=teal!35!black,title={Before you move on}]
What does \ru{отходить} mean? (\textquotedblleft{}To pull out, to leave (of a train); to step away\textquotedblright{}.) Say it once more.
\end{tcolorbox}
\section[\texorpdfstring{\ru{объявлять} (obyavlyát)}{объявлять (obyavlyát)}]{\ru{объявлять} (obyavlyát) — to announce}

\label{lesson:RU-C415-obyavlyat}



\begin{quote}
Before the new one: say the Russian for to arrive, then the Russian for to pull out.
\end{quote}

\subsection*{You'll want to know: \ru{объявлять}}
\textbf{\ru{объявлять}} (\emph{obyavlyát}) — \textquotedblleft{}to announce\textquotedblright{}.

\begin{quote}
\ru{На} \ru{вокзале} \ru{объявляют} \ru{посадку} \ru{на} \ru{поезд}.
\end{quote}

\begin{quote}
\emph{Na vokzále obyavlyáyut posádku na póyezd.}
\end{quote}

\begin{quote}
At the station they are announcing boarding for the train.
\end{quote}

Imperfective; the perfective is \emph{\ru{объявить}} (\emph{\ru{объявлю}, \ru{объявишь}}). \emph{\ru{Объявление}} — an announcement.

\begin{cousinweb}
\emph{\ru{объявлять}} is built from \emph{\ru{об}-} and \emph{\ru{являть}} 'to show'. EtymDB traces \emph{\ru{явить}} through Old East Slavic \emph{javiti} to Proto-Slavic \emph{aviti} 'to make visible'.
\end{cousinweb}

\begin{grammarlens}[title={What you've built}]
One more word for telling, explaining and describing at B1.
\end{grammarlens}

\subsection*{Your turn}
\begin{itemize}
\raggedright
  \item \emph{Say it:} \emph{obyavlyát}
  \item \emph{Say it:} \emph{obyavlyát}, once more
  \item \emph{Say it:} a sentence of your own with \emph{obyavlyát}
  \item \emph{Recall:} say the Russian for to arrive, then the Russian for to pull out, then say \emph{obyavlyát} again
\end{itemize}

\begin{tcolorbox}[breakable,colback=teal!4,colframe=teal!35!black,title={Before you move on}]
What does \ru{объявлять} mean? (\textquotedblleft{}To announce\textquotedblright{}.) Say it once more.
\end{tcolorbox}
\section[\texorpdfstring{\ru{успевать} (uspevát)}{успевать (uspevát)}]{\ru{успевать} (uspevát) — to make it (in time), to have time (to)}

\label{lesson:RU-C415-uspevat}



\begin{quote}
Before the new one: say the Russian for to pull out, then the Russian for to announce.
\end{quote}

\subsection*{You'll want to know: \ru{успевать}}
\textbf{\ru{успевать}} (\emph{uspevát}) — \textquotedblleft{}to make it (in time), to have time (to)\textquotedblright{}.

\begin{quote}
\ru{Мы} \ru{бежим} \ru{и} \ru{всё}-\ru{таки} \ru{успеваем} \ru{на} \ru{поезд}.
\end{quote}

\begin{quote}
\emph{My bezhím i vsyo-taki uspeváyem na póyezd.}
\end{quote}

\begin{quote}
We run and still make the train.
\end{quote}

Imperfective; perfective \emph{\ru{успеть}} (\emph{\ru{успею}, \ru{успеешь}}). \emph{\ru{Успеть} \ru{на} \ru{поезд}} — to catch the train.

\begin{cousinweb}
\emph{\ru{успевать}} is built from \emph{\ru{у}-} and an old verb \emph{\ru{спеть}} 'to get on, to succeed', the root also found in \emph{\ru{спешить}} 'to hurry'. EtymDB lists \emph{\ru{успевать}} without a source; the link with \emph{\ru{спешить}} is the usual one.
\end{cousinweb}

\begin{grammarlens}[title={What you've built}]
One more word for telling, explaining and describing at B1.
\end{grammarlens}

\subsection*{Your turn}
\begin{itemize}
\raggedright
  \item \emph{Say it:} \emph{uspevát}
  \item \emph{Say it:} \emph{uspevát}, once more
  \item \emph{Say it:} a sentence of your own with \emph{uspevát}
  \item \emph{Recall:} say the Russian for to pull out, then the Russian for to announce, then say \emph{uspevát} again
\end{itemize}

\begin{tcolorbox}[breakable,colback=teal!4,colframe=teal!35!black,title={Before you move on}]
What does \ru{успевать} mean? (\textquotedblleft{}To make it (in time), to have time (to)\textquotedblright{}.) Say it once more.
\end{tcolorbox}
